% ============================================================================
% TEMPLATE BASE PARA APRESENTAÇÕES EM BEAMER
% Bioinformática para Biologia de Sistemas
% ============================================================================

\documentclass[12pt, aspectratio=169]{beamer}

% ============================================================================
% PACOTES ESSENCIAIS
% ============================================================================
\usepackage[utf8]{inputenc}
\usepackage[T1]{fontenc}
\usepackage[brazilian]{babel}
\usepackage{amsmath, amssymb, amsthm}
\usepackage{graphicx}
\usepackage{xcolor}
\usepackage{tikz}
\usetikzlibrary{arrows.meta, shapes, positioning, calc, decorations.pathreplacing, decorations.shapes, decorations.pathmorphing}
\usepackage{listings}
\usepackage{booktabs}
\usepackage{multirow}
\usepackage{hyperref}
\usepackage{chemfig}
\usepackage{chemarr}

% ============================================================================
% CONFIGURAÇÃO DO TEMA
% ============================================================================
\usetheme{Madrid}
\usecolortheme{default}

% Cores personalizadas
\definecolor{azulescuro}{RGB}{0, 51, 102}
\definecolor{azulclaro}{RGB}{51, 153, 255}
\definecolor{verdeacido}{RGB}{102, 204, 0}
\definecolor{laranjacel}{RGB}{255, 153, 51}
\definecolor{roxodna}{RGB}{153, 51, 255}

\setbeamercolor{structure}{fg=azulescuro}
\setbeamercolor{palette primary}{bg=azulescuro,fg=white}
\setbeamercolor{palette secondary}{bg=azulclaro,fg=white}
\setbeamercolor{palette tertiary}{bg=azulescuro,fg=white}
\setbeamercolor{palette quaternary}{bg=azulclaro,fg=white}
\setbeamercolor{block title}{bg=azulescuro,fg=white}
\setbeamercolor{block body}{bg=azulclaro!10,fg=black}
\setbeamercolor{block title alerted}{bg=red!80,fg=white}
\setbeamercolor{block body alerted}{bg=red!10,fg=black}
\setbeamercolor{block title example}{bg=verdeacido!80,fg=white}
\setbeamercolor{block body example}{bg=verdeacido!10,fg=black}

% ============================================================================
% CONFIGURAÇÃO DE CÓDIGO
% ============================================================================
\lstset{
    language=Python,
    basicstyle=\ttfamily\small,
    keywordstyle=\color{azulescuro}\bfseries,
    commentstyle=\color{verdeacido},
    stringstyle=\color{laranjacel},
    numbers=left,
    numberstyle=\tiny\color{gray},
    stepnumber=1,
    numbersep=5pt,
    backgroundcolor=\color{white},
    showspaces=false,
    showstringspaces=false,
    showtabs=false,
    frame=single,
    rulecolor=\color{black},
    tabsize=2,
    captionpos=b,
    breaklines=true,
    breakatwhitespace=false,
    escapeinside={\%*}{*)},
    xleftmargin=10pt,
    xrightmargin=10pt
}

% Definir estilo para R
\lstdefinestyle{Rstyle}{
    language=R,
    basicstyle=\ttfamily\small,
    keywordstyle=\color{azulescuro}\bfseries,
    commentstyle=\color{verdeacido},
    stringstyle=\color{laranjacel}
}

% Definir estilo para MATLAB
\lstdefinestyle{MATLABstyle}{
    language=Matlab,
    basicstyle=\ttfamily\small,
    keywordstyle=\color{azulescuro}\bfseries,
    commentstyle=\color{verdeacido},
    stringstyle=\color{laranjacel}
}

% ============================================================================
% COMANDOS PERSONALIZADOS
% ============================================================================

% Comando para equações destacadas
\newcommand{\destaque}[1]{%
    \begin{alertblock}{}
        \begin{center}
            #1
        \end{center}
    \end{alertblock}
}

% Comando para definições
\newcommand{\definicao}[2]{%
    \begin{block}{Definição: #1}
        #2
    \end{block}
}

% Comando para exemplos
\newcommand{\exemplo}[2]{%
    \begin{exampleblock}{Exemplo: #1}
        #2
    \end{exampleblock}
}

% Comando para observações importantes
\newcommand{\importante}[1]{%
    \begin{alertblock}{Importante!}
        #1
    \end{alertblock}
}

% Comandos matemáticos úteis
\newcommand{\R}{\mathbb{R}}
\newcommand{\N}{\mathbb{N}}
\newcommand{\Z}{\mathbb{Z}}
\newcommand{\dif}{\mathrm{d}}
\newcommand{\parcial}[2]{\frac{\partial #1}{\partial #2}}
\newcommand{\derivada}[2]{\frac{\dif #1}{\dif #2}}
\newcommand{\vect}[1]{\boldsymbol{#1}}

% Comandos biológicos
\newcommand{\gene}[1]{\textit{#1}}
\newcommand{\proteina}[1]{\textsc{#1}}
\newcommand{\especie}[1]{\textit{#1}}

% ============================================================================
% CONFIGURAÇÃO DE RODAPÉ
% ============================================================================
\setbeamertemplate{footline}{
  \leavevmode%
  \hbox{%
  \begin{beamercolorbox}[wd=.33\paperwidth,ht=2.25ex,dp=1ex,center]{author in head/foot}%
    \usebeamerfont{author in head/foot}\insertshortauthor
  \end{beamercolorbox}%
  \begin{beamercolorbox}[wd=.34\paperwidth,ht=2.25ex,dp=1ex,center]{title in head/foot}%
    \usebeamerfont{title in head/foot}\insertshorttitle
  \end{beamercolorbox}%
  \begin{beamercolorbox}[wd=.33\paperwidth,ht=2.25ex,dp=1ex,right]{date in head/foot}%
    \usebeamerfont{date in head/foot}\insertshortdate{}\hspace*{2em}
    \insertframenumber{} / \inserttotalframenumber\hspace*{2ex}
  \end{beamercolorbox}}%
  \vskip0pt%
}

% ============================================================================
% CONFIGURAÇÃO DE NAVEGAÇÃO
% ============================================================================
\setbeamertemplate{navigation symbols}{}

% ============================================================================
% CONFIGURAÇÃO DE ITEMIZE/ENUMERATE
% ============================================================================
\setbeamertemplate{itemize items}[circle]
\setbeamertemplate{enumerate items}[default]

% ============================================================================
% FIM DO TEMPLATE
% ============================================================================
